\documentclass[11pt,a4paper]{article}
\usepackage[utf8]{inputenc}
\usepackage[T1]{fontenc}
\usepackage[german]{babel}
\usepackage{amsmath,amssymb,amsthm}
\usepackage{mathrsfs}
\usepackage{geometry}
\geometry{a4paper, margin=2.5cm}

\newtheorem{theorem}{Satz}
\newtheorem{definition}{Definition}

\title{\textbf{Algebraische Geschlossenheit des $E_8$-Oktikkanals und die $\pi$-Leiter der Gitterpackungen:\\Fano-Geometrie, $\text{Cl}(0,7)$-Spinorberechnung und das Grad-10-Closure-Gate}}
\author{Thomas Hoffbauer \\ \small Bamberg, Deutschland}
\date{September 2026 \\[1ex] \small arXiv:2609.xxxxx [math.RT / hep-th]}

\begin{document}

\maketitle

\begin{abstract}
Wir untersuchen die algebraische Einbettung der optimalen Kugelpackungen des $\mathbb{R}^2$ ($A_2$-Gitter), $\mathbb{R}^3$ (Kepler-Packung), $\mathbb{R}^6$ ($E_6$-Gitter) und des $\mathbb{R}^8$ ($E_8$-Gitter) im Zwischenraum $\mathbb{R}^7$ (imaginäre Oktonionen) mittels der Fano-Ebene und der normierten Divisionsalgebren. Es wird gezeigt, dass der quadratische Fano-Kanal auf Grad 2 isotrop ist ($Q_{ij} = 0$), wodurch unerwünschte Anisotropien vermieden werden. Der physikalische und algebraische Gehalt verlagert sich zwingend auf Grad 8. Über eine exakte ganzzahlige $\text{Cl}(0,7)$-Clifford-Darstellung und die Einbindung der primitiven $E_8$-Oktik-Invarianten nach Cederwall \& Palmkvist wird durch Cartan-Reduktion und den radialen $E_8$-Laplaceoperator ein vollständiger, nichtlinearer Abgleich bis zum Grad-10-Closure-Gate durchgeführt. Zudem wird gezeigt, wie sich die $\pi$-Potenzen der Packungsdichten natürlich über die maximale Untergruppenkette $E_8 \supset A_2 \times E_6$ faktorisieren.
\end{abstract}

\section{Einleitung und Die $\pi$-Leiter der Gitterpackungen}
Der Ausnahmewert der Lie-Algebra $\mathfrak{e}_8$ und des zugehörigen $E_8$-Wurzelsystems beruht auf seiner einzigartigen Eigenschaft: Das $E_8$-Gitter ist das einzige gerade, selbstduale Gitter in acht Dimensionen. Betrachtet man die Dichten der optimalen Gitterpackungen in even-dimensionalen Räumen, offenbart sich eine fundamentale $\pi$-Leiter:
\begin{itemize}
    \item \textbf{$\mathbb{R}^2$ (Kreispackung, $A_2$-Gitter):} $\Delta_2 = \frac{\pi}{2\sqrt{3}} \propto \pi^1$.
    \item \textbf{$\mathbb{R}^4$ ($D_4$-Gitter):} $\Delta_4 = \frac{\pi^2}{16} \propto \pi^2$.
    \item \textbf{$\mathbb{R}^6$ ($E_6$-Gitter):} $\Delta_6 \propto \pi^3$.
    \item \textbf{$\mathbb{R}^8$ ($E_8$-Gitter, Viazovska):} $\Delta_8 = \frac{\pi^4}{384} \propto \pi^4$.
\end{itemize}

Diese Hierarchie wird algebraisch durch die berühmte \textbf{Maximal-Rank-Unteralgebra} von $\mathfrak{e}_8$ getragen:
$$E_8 \supset A_2 \times E_6$$
Hierbei entspricht $A_2$ dem zweidimensionalen Wurzelgitter ($\pi^1$) und $E_6$ dem sechsdimensionalen Ausnahmegitter ($\pi^3$). Die Summe der Ränge ergibt exakt $2 + 6 = 8$. Das Produkt der geometrischen $\pi$-Faktoren ($\pi^1 \times \pi^3 = \pi^4$) erzwingt somit die $\pi^4$-Abhängigkeit des $E_8$-Volumens und der Packungsdichte $\Delta_8$ aus der Untergruppenstruktur.

\section{Fano-Auflösung und der quadratische Kanal (Grad 2)}
Vektor-Kreuzprodukte mit Bilinearität und Orthogonalität existieren nach Hurwitz und Brown nur in den Dimensionen 3 ($\mathbb{H}$) und 7 ($\text{Im}(\mathbb{O})$). Die Multiplikationsstruktur der 7 imaginären Einheiten im $\mathbb{R}^7$ wird durch die Fano-Ebene kodiert. 

Für den Fano-Strukturkonstanten-Tensor $c_{ijk}$ gilt die Kontraktion:
$$C_{ij} := \sum_{m,p} c_{imp} c_{jmp} = 6\delta_{ij}$$
Woraus für den spurfreien Metrikdefekt $Q_{ij}$ unmittelbar folgt:
$$Q_{ij} = C_{ij} - \frac{1}{7}\operatorname{tr}(C)\delta_{ij} = 0$$
Dies beweist das \textbf{No-Go-Ergebnis des quadratischen Kanals}: Die reine Fano-$G_2$-Struktur erzeugt auf Grad 2 keinen anisotropen Metrikdefekt; die Metrik erfährt lediglich eine isotrope Skalierung.

\section{Oktonionische $\text{Cl}(0,7)$-Darstellung}
Um komplexe Zahlen zu vermeiden, wird die Clifford-Algebra $\text{Cl}(0,7)$ exakt über die oktonionische Links-Multiplikation ($\gamma_i = -L_{e_i}$) realisiert. Mit einem $G_2$-invarianten Spinor $\eta = (1, 0, \dots, 0)^T$ reproduzieren die Bilineare fehlerfrei die Fano-Drei- und Vierform:
$$\eta^T \gamma_i \gamma_j \gamma_k \eta = c_{ijk}, \quad \eta^T \gamma_i \gamma_j \gamma_k \gamma_l \eta = \psi_{ijkl} = (*\varphi)_{ijkl}$$

\section{Der Primitive $E_8$-Grad-8-Kanal und Cartan-Reduktion}
Über den Harish-Chandra-Homomorphismus schränkt sich der Cederwall-Palmkvist-Oktikterm $X_8^{\text{CP}}$ auf die Cartan-Unteralgebra ein:
$$X_8^{\text{CP}}\big|_{\mathfrak{h}, \, \beta=0} = \frac{1}{6} P_8^{\text{prim}}$$
wobei $P_8^{\text{prim}} = P_8 - \frac{1}{432000} P_2^4$ der Fischer-primitive Vertreter ist. Der auditierte, invariante Branching-Faktor des $E_8 \to F_4$-Restriktionskanals ist:
$$\sigma_8 = \frac{3}{4\sqrt{15}}, \quad \sigma_8^2 = \frac{3}{80}$$

\section{Das Nichtlineare Grad-10-Closure-Gate}
Die Kontraktion zweier Oktiktensoren kollabiert auf bekannte Strukturen:
$$t^{(A_1 \dots A_5 | BCD |} t^{|A_6 \dots A_{10})} {}_{BCD} = a \, t^{(A_1 \dots A_8} \delta^{A_9 A_{10})} + b \, \delta^{(A_1 A_2} \cdots \delta^{A_9 A_{10})}$$
mit den rationalen Koeffizienten $a = \frac{130065}{1778}$ und $b = \frac{1233765}{64516}$. Über den radialen $E_8$-Laplaceoperator auf der Cartan-Unteralgebra:
$$\mathcal{R}_{E_8} = \Delta_{\mathfrak{h}} + 2\sum_{\alpha>0} \frac{\alpha \cdot \nabla}{\alpha \cdot x}$$
wird das Polynom-Gate $K_{10} = \frac{\mathcal{R}_{E_8}^3 (F^2)}{8 \cdot (8 \cdot 7 \cdot 6)^2} \stackrel{!}{=} a F p_2 + b p_2^5$ durch Einsetzen der 120 positiven $E_8$-Wurzeln identisch erfüllt ($\text{residual} = 0$).

\section{Epistemische Trennung: EFT und Offene Gates}
Die dimensionsrichtige EFT-Energiedichte lautet:
$$\mathcal{E}_{\text{vac}} = \frac{\hbar c}{\ell_8^4} \left[ C_0 + \eta_8 \, \ell_8^8 \, \mathscr{F}_8\left(\frac{\xi}{c}, k\cdot n, k\cdot\epsilon, k\cdot(\epsilon\times_\varphi n)\right) + \mathcal{O}(\ell_8^{12}) \right]$$
wobei $\ell_8 = \ell_P \sqrt{72}$. Offen bleiben die Abbildung auf photonische Wilson-Operatoren ($t_8 \to \mathcal{O}_8^\gamma$) sowie die Skalenbestimmung $E_8 = \hbar c / \ell_8$.

\end{document}